\section{Discussion}
\label{sec:discussion}

\subsection{Principal findings}

\emph{In one sentence:} four findings --- degeneracy governs everything; both sources gain over S1 out of sample and remain inseparable; the error-signal claim survives only in weakened, model-dependent form; fragility tracks answer disagreement.

The study yields four principal findings. First, the effectiveness of every test-time aggregation strategy is governed by a precondition that the literature rarely states explicitly: the model must produce diverse answers at all. For the flagship models in our pool, five of six benchmarks sit in a degenerate regime in which six stochastic samples and six paraphrases return the same answer for almost every question, and in that regime self-consistency, paraphrase ensembling, and the combined ensemble reproduce the single-call baselines exactly, with the hybrid allocation differing by at most one question. This part of the degeneracy result is structural rather than empirical: when all votes coincide, every aggregation of them returns the same answer by construction. The empirical content lies in how common the coincidence is, and in its consequence that agreement-based uncertainty scores are uninformative exactly there.

Second, measured against the correct $k{=}1$ baseline, the single-sample expectation S1, the two diversity sources deliver small, reliable, and statistically inseparable aggregation gains on the held-out suite: pooled over its four fully covered models (476 model--question pairs), self-consistency gains $+2.4$ points over S1, the paraphrase ensemble $+4.9$, and the twelve-vote combination $+4.1$, all with intervals excluding zero and surviving Holm correction within the designated primary family (Section~\ref{sec:stats}); the paraphrase margin over self-consistency ($+2.5$, CI $-0.6$ to $+5.7$) does not separate from zero, and after removing the paraphrase difficulty margin ($+1.9$ at the single-call level) the like-for-like aggregation comparison is $+3.0$ against $+2.4$ --- roughly a factor of 1.25, not the factor of two that the raw headline suggests. On the development suite, whose per-model coverage was truncated unevenly (12 to 90 questions), the same pooled comparison sits at noise (SC-6 $+0.9$, PE-6 $-0.9$ over 112 pairs), and the per-model gains measured against greedy decoding are largely artifacts of a greedy reference that diverges from the sampled expectation on specific endpoints; we therefore treat the held-out suite as the primary evidence and the development suite as the design phase whose role was to fix the conditions, strata, and rules. Where the modal answer is wrong, the two sources pull in opposite directions, and the sign of the margin is model- and dataset-conditional: the paraphrase-side margin appears where the canonical instruction under-elicits --- on quantitative tasks, where the instruction-variant analysis of Section~\ref{sec:heldoutmech} shows the best single variant (selected post hoc) already exceeding the six-variant ensemble --- and reverses where paraphrases shift verbal semantics (LogiQA for one model, PIQA for another); the allocation that realizes the gain varies across models. Because every paraphrase vote mixes prompt variation with decoding noise and no greedy-paraphrase control exists on the held-out suite, these margins characterize paraphrase ensembles with sampling noise, not pure prompt-space effects.

Third, paraphrase agreement functions as an error signal in its own right, and its error coverage overlaps only partially with that of sample agreement --- but the per-model regularity we first read from the development suite does not survive. The development suite could support its readings only descriptively (its informative models carried one to three ensemble errors, below the five-error rule of Section~\ref{sec:method}); the held-out suite recomputes them with twelve to twenty-one errors per model. What survives: of the 72 pooled errors, sample dissent flags 49 and paraphrase dissent 46, with fifteen caught only by the sample side and twelve only by the paraphrase side (the union's coverage is an upper bound by construction; the partial disjointness is the empirical content), and pooled AUROC is comparable across the two signals (0.75 and 0.74). What does not survive: the development-suite reading that sample agreement is decisively stronger for Agnes3.0-Flash (0.96 against 0.60) rested on three errors and is overturned by the held-out recomputation (0.80 against 0.81); on the fresh questions the stronger single signal is model-dependent, not a tracker of the model's stronger diversity source. The summed disagreement score is numerically the strongest pooled detector (AUROC 0.81 against 0.75 and 0.74); we label that ranking descriptive, since its interval overlaps the singles'. The usable signal is nevertheless mostly a regime signal: with the definitional degenerate stratum excluded, per-model AUROCs fall to between 0.49 and 0.86, so agreement separates risky-regime questions from safe ones --- a 35\% error rate where any vote dissents against 3.6\% where none does, pooled --- far better than it ranks risk within the risky set.

Fourth, fragility, the fraction of questions whose answer flips under paraphrasing alone, tracks answer disagreement, concentrates in stem-locus paraphrasing (17.4\% pooled flips on the held-out suite against 9.6\% under instruction variation, with the stem-locus flips again centred on LogiQA), and appears in the same model-task cells that carry disagreement, which supports the interpretation that paraphrase disagreement and sampling disagreement are two observable symptoms of the same underlying instability. The development suite's stronger claim --- that instruction-locus flips never occur --- proved to be a property of its question draw rather than a structural law: on the held-out set, instruction-locus flips appear wherever the instruction variants differ in elicitation quality, almost entirely on SAT-Math, which is also where the instruction-variant analysis of Section~\ref{sec:heldoutmech} locates a single-variant effect larger than the aggregation margins.

\subsection{A practical allocation rule, and how it performs}

\emph{In one sentence:} a probe-and-route rule saves 13--21 percent of the calls on fresh questions at matched or slightly lower accuracy, and its value is the call saving, not accuracy; a simpler unanimity rule preserves self-consistency at two to three calls.

The regime structure suggests a simple procedure for practitioners with a fixed per-question budget. Probe the question first with a small number of calls, compute sample agreement and paraphrase agreement, and branch on the result. If both agreement scores equal one, additional calls of either kind are unlikely to change the answer, and the budget is better spent elsewhere; the degenerate regime is identifiable from a handful of calls at negligible cost. If sample agreement is high but paraphrase agreement is low, the model is prompt-brittle, and the remaining budget should go to paraphrases. If the reverse holds, the budget should go to resampling. If both scores are low, split the remaining budget between the two sources. The rule was evaluated by simulation on the development matrices (a consistency check, not a validation, since it was designed on the same data; \ref{app:dev}) and transferred to the held-out suite in Section~\ref{sec:heldoutrule}: there it saves 13 to 21 percent of the six-call budget, exceeds every fixed allocation on one model by a single question, loses one to two questions to the best fixed conditions on two models, and sits seven questions below the paraphrase vote on Agnes3.0 --- routing toward paraphrases when answer-space diversity is low remains the regime where the most is gained, and it is exactly the regime the probe's two-call answer-space check undersamples. An even simpler sequential rule that draws canonical samples until the first unanimity and otherwise falls back to the six-sample majority preserves SC-6 accuracy within one question everywhere at a mean budget of two to three calls; this is a special case of the early-stopping self-consistency family \citep{aggarwal2023lets,li2024escape}, and our simulation reproduces that family's finding in the direct-answering, near-degenerate regime. Because agreement statistics are cheap to compute online, the rules add no model training and no human labeling, and both simulations are reproducible from the logs released with the submission.

\subsection{Implications for evaluation and for uncertainty estimation}

\emph{In one sentence:} reported scores should state the aggregation configuration and the single-sample expectation S1 alongside greedy; sampling-based confidence measures must be validated per regime.

Two methodological consequences follow. For benchmark evaluation, our results add an aggregation dimension to the multi-prompt argument of \citet{mizrahi2024state}: the same question answered under one prompt with one decode, under six paraphrases, or under twelve combined calls yields systematically different accuracies in the high-entropy regime, so reported scores should state the aggregation configuration, and comparisons across papers that aggregate differently are not like for like. The single-sample baseline S1 deserves the same reporting discipline, because greedy decoding and a single stochastic draw differ on a minority of model--endpoint configurations --- absent on three of the four held-out endpoints, present on one --- and are statistically indistinguishable on the rest; the divergence is an endpoint property to be checked per deployment, not a universal offset. For uncertainty estimation, the degenerate regime is a warning whose scope we state precisely. Our measurements cover multiple-choice questions answered directly, where sampling-based confidence measures, including simple agreement heuristics, are uninformative when the sample distribution collapses, and our data show that this collapse is the common case for strong models on saturated tasks. Semantic-entropy methods for free-form generation \citep{kuhn2023semantic,farquhar2024detecting} face an analogous risk when the sampled paraphrase classes collapse, but our design does not test that setting, and the caution transfers by argument rather than by measurement. Paraphrase probes partially restore the signal because they perturb the input rather than the decoder, but they too inherit the collapse when the model is robust in both views, as our flagship models are on five of six benchmarks. Confidence methods should therefore be validated per regime rather than averaged over mixed test sets, and deployments should verify that the answer distribution carries variance before trusting agreement-based scores.

\subsection{Limitations}
\label{sec:limits}

\emph{In one sentence:} nine limitations bound the scope of the conclusions, stated without defence.

We list them plainly. First, sample size and power: fifteen questions per benchmark and completed sets of 12 to 90 questions per model support paired within-model comparisons of large effects, but the cluster-bootstrap intervals are wide, TOST does not establish equivalence for most pairs, and differences of a few points between strategies would go undetected; the 119-question held-out set multiplies per-model coverage for four systems, while two of its six systems remain partially covered, and scaling each benchmark to hundreds of questions remains the single change most likely to sharpen the conclusions. Second, provider rate limits and free-tier quotas constrained coverage unevenly across models; free-tier access to a given model--endpoint configuration can be withdrawn at any time, coverage follows a fixed question order rather than adaptive selection, but cross-model aggregates weight models unequally. Third, the models are commercial systems accessed through third-party endpoints, one of which is an aggregator rather than a first-party API: we can neither verify that a relay serves the requested model identity and version against the provider's own API nor confirm that the requested temperature is honored, especially for reasoning-style models; no condition in this study was replicated on a first-party API or an open-weight model, so all findings describe the accessed endpoints as of October 2026 and need not transfer to other versions or serving configurations. We release per-call timestamps, token usage, and raw outputs so that the configuration is auditable; an official-API or open-weight replication remains the cleanest resolution and is left open. Fourth, paraphrase equivalence is certified only by automatic filters --- embedding similarity with a fixed threshold plus a bidirectional LLM entailment check --- and semantic equivalence of adversarial questions is not fully formalizable; the filters bound equivalence in distribution, not item by item, and a human equivalence spot-check has not yet been run\TODO{planned human equivalence spot-check on LogiQA stem rewrites: 100-item annotation table and summary script provided with the submission, result pending}. A related caveat is that Qwen3.7-Flash, the stem-paraphrase generator, is itself evaluated in the held-out suite, so its paraphrase-condition cells carry a self-preference risk, and its partial coverage (60 of 119) keeps them out of the pooled analyses. Fifth, the paraphrase side mixes two perturbations that the design cannot separate out of sample: every paraphrase vote varies the prompt \emph{and} carries $\tau = 0.7$ decoding noise, and the greedy-paraphrase control that isolates the prompt component was collected on the development suite only, so all held-out paraphrase margins are margins of paraphrase ensembles with sampling noise, not pure prompt-space effects. Sixth, the SAT-Math instruction-variant analysis is post hoc --- the best variant was selected on the same data it scores, so its 95.0\% is an optimistic ceiling --- and the control that would test whether the canonical instruction is suboptimal under aggregation, self-consistency under the pre-selected best instruction, has not been run\TODO{pre-registered follow-up experiment, script provided, result pending (Section~\ref{sec:heldoutmech})}. Seventh, the multiple-choice tasks use direct answering while GSM8K uses chain of thought; the aggregation comparison is internally valid within each answering mode, and the follow-up condition of \ref{app:cot} measures the mode effect directly, but absolute accuracies are not comparable across modes. Eighth, the high-dissent regime is sampled through fewer model-task cells than the degenerate regime, resting primarily on logical reasoning and selected quantitative tasks; the MMLU-Pro probe of \ref{app:mmlupro} begins to broaden this, at small scale. Ninth, error-detection metrics are computed against ensemble references that share votes with the detectors, and the adaptive allocation rule was designed and evaluated on development-suite matrices before its held-out transfer, so both carry a same-data caveat that the held-out set mitigates but does not eliminate; an error reference from fully independent draws would still be cleaner. A tenth, cross-cutting point: all conditions use a fixed budget of six or twelve calls with majority voting; ranked voting, weighted aggregation, or learned vote combination could change the budget-matched comparison.

\subsection{Future work}

Three extensions follow naturally. The adaptive allocation rule can be moved online, interleaving probing and allocation per question, and its threshold can be tuned against a target budget rather than fixed at unanimity. The comparison can be extended to free-form generation, where semantic clustering replaces majority voting and paraphrase ensembles interact with surface-form competition effects \citep{holtzman2021surface}, and where the degeneracy phenomenon has direct consequences for semantic-entropy hallucination detection that our multiple-choice design cannot measure. Finally, the fragility statistic can serve as a benchmark diagnostic in its own right, complementing accuracy with a measure of how strongly a model's answers depend on the surface form of the question.
